There are some systems that do not use indexes, or very sparingly.
Instead, they keep metadata about the blocks and store them in memory and use it to check if the data in the block is needed.

This is common in column-store engines such as DuckDB and Snowflake.
This metadata can be as simple as min/max for each column,
or complete stats, bitmaps, precomputed aggregates, or even bloom filters.
